\section{Matrix balls and sparse completion cones}
\label{sec:spectral-chordal}

This section compares two lifts of products of spectral-norm matrix
balls in which one cone factor serves several blocks.  The first places
several blocks under one scale variable, which reduces the number
of cone coordinates and the ambient barrier parameter.  The second uses a
positive-semidefinite completion cone, which keeps a sparse matrix
structure.  After the scale or the diagonal blocks are fixed, both
constructions have exactly the same barrier on the product of matrix balls.
The spectral and completion barriers used below are classical
\citep{CKV2022natural,ADV2013}.  Our results are factorizations of the
full slack rows $1-\ip{X_a}{Z}$ by these cones, sharp bounds on cone and
face dimensions, and the equality of the two reduced Newton systems after
this fixing.

\subsection{Mixed contact rank with split rows}

For $\mathbb F\in\{\R,\C,\HH\}$ let
$\delta=\dim_{\R}\mathbb F$ and use the real Frobenius pairing
$\ip XZ=\operatorname{Re}\tr(X^*Z)$.
The extreme points of the spectral-norm ball in
$\mathbb F^{r\times c}$, $1\leq r\leq c$, are the coisometries
\[
 \mathcal E_{r,c}=\{X:XX^*=I_r\};
\]
those of its nuclear-norm polar are
$\mathcal Z_{r,c}=\{uv^*:\norm u=\norm v=1\}$.
These descriptions follow from singular-value decomposition: a contraction
with a singular value below one is a nontrivial average of contractions,
whereas equality on every row forces any such averaging at a coisometry
to be trivial.  Similarly, the nuclear unit ball is the convex hull of the
rank-one unit matrices, which are exposed by suitable contractions.

Consider $k$ blocks $X_a\in\mathcal E_{r_a,c_a}$ over fields
$\mathbb F_a$, with $\delta_a=\dim_{\R}\mathbb F_a$.  Suppose the full
slack rows $1-\ip{X_a}{Z}$ have a globally $C^1$ factorization with
dual maps labelled by the row:
\begin{equation}\label{eq:spectral-split-slack}
 1-\ip{X_a}{Z}
 =\sum_i\ip{A_i(X_1,\ldots,X_k)}{B_i^a(Z)},
 \qquad A_i\in K_i,\quad B_i^a\in K_i^*.
\end{equation}
Here the nonray proper cones have real dimensions $m_i\geq3$;
ray factors contribute no mixed rank.  Put
\[
 q_i=m_i-2,\quad
 f_i=\max_{0\ne y\in K_i^*}\dim(K_i\cap y^\perp),\quad
 \bar f_i=\min(f_i,k),\quad \kappa_a=\delta_a c_a-1.
\]
All dimensions are real.

\begin{proposition}[Rank and face bounds at spectral contacts]\label{prop:spectral-contact}
At every contact $\ip XZ=1$, the mixed pairing
$(H,\dot Z)\mapsto-\ip H{\dot Z}$ on
$T_X\mathcal E_{r,c}\times T_Z\mathcal Z_{r,c}$ has rank
$\delta c-1$.  At a simultaneous contact
$\ip{X_a}{Z_a}=1$ for all $a$, let
\[
 t_i^a=\rank\bigl[(H,\dot Z)\mapsto
       -\ip{D_aA_i[H]}{DB_i^a[\dot Z]}\bigr],
 \qquad I_i=\{a:t_i^a>0\},
\]
and call $I_i$ the rows active in factor $i$.  Then
\begin{align}
 t_i^a&\leq q_i,& \sum_i t_i^a&\geq\kappa_a,
       & |I_i|&\leq\bar f_i,\label{eq:spectral-channel-budgets}\\
 \sum_{a\in I_i}t_i^a&\leq q_i,&
 1+\sum_{b\in I_i\setminus\{a\}}t_i^b
   &\leq\dim(K_i\cap B_i^a(Z_a)^\perp)
       &&(a\in I_i).\label{eq:spectral-face-budget}
\end{align}
If $m_i\leq d$ for an integer $d\geq3$, these imply
\begin{equation}\label{eq:spectral-summed-budgets}
 \sum_i\bar f_i\geq
       \sum_a\left\lceil\frac{\kappa_a}{d-2}\right\rceil,
 \qquad
 \sum_iq_i\geq\sum_a\kappa_a.
\end{equation}
In particular, if the $L$ nonray factors satisfy
$f_i\leq f$ for a common $f\geq1$, then
\[
 L\geq\max\left\{
       \left\lceil\frac1f\sum_a
          \left\lceil\frac{\kappa_a}{d-2}\right\rceil\right\rceil,
       \left\lceil\frac{\sum_a\kappa_a}{d-2}\right\rceil
       \right\}.
\]
\end{proposition}
\begin{proof}
Contact implies $Xv=u$ and $X^*u=v$ by equality in the contraction
inequality.  Unitary changes of bases therefore give
$X=(I_r\;0)$ and $u=v=e_1$.  Write a coisometry tangent as
$H=(A\;B)$ with $A^*=-A$.
Since $(u,v)$ and $(uq,vq)$ give the same $uv^*$ for a unit scalar
$q\in\mathbb F$, we may assume $u^*\dot u=0$.  The pairing
$\ip H{\dot u e_1^*+e_1\dot v^*}$ splits into independent parts: the
entries $A_{j1}$,
$j>1$, paired with the off-axis differences of $\dot u$ and $\dot v$;
the first row of $B$, paired with the tail of $\dot v$; and the imaginary
entry $A_{11}$, paired with the imaginary entry $\dot v_1$.
They have respectively $\delta(r-1)$, $\delta(c-r)$, and $\delta-1$
real dimensions, and the pairing on each part is nondegenerate.
Their sum is $\delta c-1$.  This argument includes real scalar rows,
where the rank can be zero.  In the quaternionic case all pairings are
real Euclidean pairings, so cyclic rearrangements of quaternionic traces
are not required.

The face argument does not need nondegeneracy of the full
spectral pairing.  At contact, each summand in
\eqref{eq:spectral-split-slack} vanishes.  If either factor is zero,
its derivative is zero because a differentiable map into a pointed cone
has zero derivative at the vertex.  Otherwise
Lemma~\ref{lem:universal-channel} bounds its mixed rank by $m_i-2$.  Mixed differentiation and rank
subadditivity give $\sum_i t_i^a\geq\kappa_a$.

Fix $i$ and an active row $a$.  On the contact cylinder, where $X_a$
stays in contact with $Z_a$ and the other blocks vary,
nonnegativity gives
\[
 \ip{D_bA_i[H_b]}{DB_i^a[\dot Z]}=0\quad(b\ne a),
 \qquad \ip{A_i}{DB_i^a[\dot Z]}=0.
\]
Indeed, at every point of that cylinder the summand, as a function of
$Z$, has minimum zero at $Z_a$; its vanishing $Z$-derivative is the
second identity, and differentiating it along $X_b$ gives the first.
Choose $E_a\subseteq\operatorname{im}D_aA_i$ of dimension $t_i^a$
on which pairing with $DB_i^a$ is injective.  These identities show that
$\R A_i\oplus\bigoplus_{a\in I_i}E_a$ is a direct sum. For any
fixed active row $a_0$, the nonzero vector $b=B_i^{a_0}(Z_{a_0})$
annihilates all of $\operatorname{im}DA_i$: the nonnegative function
$X'\mapsto\ip{A_i(X')}{b}$ has value zero at the contact $X$.
It also annihilates $A_i$, so the direct sum lies in $b^\perp$ and
$1+\sum_a t_i^a\leq m_i-1$. If there is no active row, the total
rank bound is immediate. Moreover,
$A_i$ and every $E_b$ with $b\ne a$ lie in the span of the exposed face
$K_i\cap B_i^a(Z_a)^\perp$.  This proves
\eqref{eq:spectral-face-budget} and $|I_i|\leq f_i$.
Finally, each row $a$ is active in at least
$\lceil\kappa_a/(d-2)\rceil$ factors. Counting these incidences and summing the mixed-rank
bounds proves \eqref{eq:spectral-summed-budgets}. Since
$\sum_iq_i\leq(d-2)L$, the count bound follows.
\end{proof}

These are local necessary conditions for split rows; they include no
global topological bound, and we do not claim that every bound is
attained.  The mixed rank is $\delta c-1$, not
the dimension of the polar extreme manifold: some polar directions move
within the contact set and do not enter the mixed pairing.

\subsection{Whole-row grouping and exact spectral parameters}

We return to the whole-row model of Section~\ref{sec:whole-rows}.
Here take $\mathbb F_a\in\{\R,\C\}$, with
$\delta_a=\dim_{\R}\mathbb F_a$, and orient each matrix with $r_a\leq c_a$.
For a nonempty group $G$ define
\[
 \mathcal S_G=\{(t,(X_a)_{a\in G}):\norm{X_a}_{\rm op}\leq t\},
 \quad N_G=\sum_{a\in G}\delta_a r_ac_a,
 \quad R_G=\sum_{a\in G}r_a,
 \quad \lambda_a=\delta_a(r_a+c_a-1).
\]
The dual cone is
$\{(\alpha,(Z_a)):\alpha\geq\sum_a\norm{Z_a}_*\}$, under the real
pairing.  Thus the polynomial maps
\[
 A_G(X)=(1,(X_a)_{a\in G}),\qquad
 B_G^a(Z)=(1,(0,\ldots,-Z,\ldots,0))
\]
represent the entire row $1-\ip{X_a}{Z}$.

\begin{corollary}[Sharp whole-row spectral caps]\label{cor:spectral-caps}
In the whole-row model, a group $G$ can be carried by one proper
cone of dimension at most $D$ and maximal proper exposed-face dimension
at most $F$ if and only if
\begin{equation}\label{eq:spectral-group-caps}
 N_G\leq D-1,\qquad N_G-\min_{a\in G}\lambda_a\leq F-1.
\end{equation}
The minimum factor count is the minimum number of groups in a partition
of the rows satisfying these inequalities; the cones $\mathcal S_G$ attain it.
For $k$ identical $r\times c$ rows over a field of real dimension $\delta$,
one group carries at most
\[
 Q=\min\left\{
    \left\lfloor\frac{D-1}{\delta rc}\right\rfloor,
    \left\lfloor\frac{F+\delta(r+c-1)-1}{\delta rc}\right\rfloor
    \right\}.
\]
If $Q\geq1$, the exact count is $\lceil k/Q\rceil$.
\end{corollary}
\begin{proof}
For a rank-one polar support, contact fixes a $1\times1$ corner of
$X_a$ and sets both off-diagonal blocks to zero; the remaining contraction has
$(r_a-1)(c_a-1)$ entries over $\mathbb F_a$.  Its affine face dimension
is therefore $\delta_a(r_a-1)(c_a-1)$, with codimension $\lambda_a$.
Theorem~\ref{thm:whole-row-caps} gives necessity.
For sufficiency, the dimension of $\mathcal S_G$ is $1+N_G$.
More generally, a dual boundary support with nonzero blocks $J$ of
ranks $h_a$ exposes a face of dimension
\[
 1+\sum_{a\notin J}\delta_a r_ac_a
   +\sum_{a\in J}\delta_a(r_a-h_a)(c_a-h_a).
\]
To see this, equality in spectral/nuclear duality fixes the supported
singular-vector corner to $-tI_{h_a}$, sets the off-diagonal
blocks to zero, and leaves an arbitrary contraction in the complementary rectangle.
The scale contributes one dimension.  Each supported block loses
$\delta_a h_a(r_a+c_a-h_a)$ coordinates; this is minimized at $h_a=1$.
Consequently the maximal proper exposed-face dimension is
$1+N_G-\min_a\lambda_a$, proving sufficiency and the count formula.
\end{proof}

For rows of different sizes, computing the minimum partition is
strongly NP-hard: for real vector rows ($r_a=1$) and a nonbinding face
cap $F\geq D-1$, it is bin packing with item sizes $c_a$ and bin
capacity $D-1$.  Since $3$-PARTITION is NP-hard with polynomially
bounded numbers, the matrix dimensions stay polynomial.  This hardness
concerns whole-row grouping, not split-row factorizations.

\begin{theorem}[Shared spectral cone and slice]\label{thm:shared-spectral-parameter}
For a partition $\mathcal G$ into $g$ nonempty groups, put
$R=\sum_a r_a$ and $N=\sum_a\delta_a r_ac_a$.
The product of the shared cones has real dimension $N+g$ and
\begin{equation}\label{eq:spectral-exact-ledger}
 \vartheta_{\rm opt}\left(\intr\prod_{G\in\mathcal G}\mathcal S_G\right)
       =R+g,\qquad
 \vartheta_{\rm opt}\left(\prod_a\{X_a:\norm{X_a}_{\rm op}<1\}\right)
       =R.
\end{equation}
Both lower bounds hold for arbitrary, possibly coupled, self-concordant
barriers.
The attaining barriers are respectively the sum over groups of
\begin{equation}\label{eq:shared-spectral-barrier}
 F_G(t,X)=-\sum_{a\in G}\log\det(t^2I_{r_a}-X_aX_a^*)
                  +(R_G-1)\log t
\end{equation}
and its fixed-scale restriction
\begin{equation}\label{eq:spectral-slice-barrier}
 \Phi(X)=-\sum_a\log\det(I_{r_a}-X_aX_a^*).
\end{equation}
\end{theorem}
\begin{proof}
The classical spectral-cone barrier is
$-\log t-\log\det(tI-XX^*/t)$, with parameter one plus the number
of rows \citep[Section 4.2]{CKV2022natural}; the same formula for complex matrices
is recorded in \citet[Section 2.6.1]{Coey2022thesis}.
Apply it to the block-diagonal matrix with blocks $X_a$; real blocks
can be embedded in a complex matrix without changing singular values.
This proves self-concordance of \eqref{eq:shared-spectral-barrier} and
its degree $R_G+1$ of logarithmic homogeneity.

For the lower bound we use Nesterov's recession bound, which does not
require logarithmic homogeneity.  For an interior point $x$, recession directions $p_j$, and positive
numbers $a_j,b_j$ satisfying
$x-\sum_j a_jp_j\in C$ and $x-b_jp_j\notin\intr C$, every
self-concordant barrier of $C$ has parameter at least $\sum_j a_j/b_j$
\citep[Theorem 3.9]{FS2023}.  We call $(x,p_j,a_j,b_j)$ a recession certificate.
Restrict the matrices of one group to their real rectangular diagonals.
Writing $r=R_G$, this gives $C_r=\{(t,\xi):t\geq\norm\xi_\infty\}$.
For $0<\epsilon<\min(1,2/r)$ choose
\[
 x=(1,(1-\epsilon)\boldsymbol1),\quad
 p_0=(1,\boldsymbol1),\quad
 p_j=(1,\boldsymbol1-2e_j)\quad(1\leq j\leq r),
\]
\[
 a_j=b_j=\epsilon/2\ (j\geq1),\quad
 a_0=1-r\epsilon/2,\quad b_0=1-\epsilon/2.
\]
All $p_j$ belong to $C_r$, $x=\sum_{j=0}^r a_jp_j$, and
$x-b_jp_j$ is on the boundary for every $j$.
The recession bound is therefore
$r+(1-r\epsilon/2)/(1-\epsilon/2)$, which tends to $r+1$.
For several groups, place each recession certificate in its own direct summand
and take $x$ to be the tuple of their interior points.  The resulting
single recession certificate gives $\sum_G(R_G+1)$ even for a coupled barrier.
These diagonal sections meet the cone interiors, so restricting any
ambient barrier to them does not increase its parameter.

For the slice, ordinary self-concordance follows by restricting the
positive-definite barrier to
$\left(\begin{smallmatrix}I_r&X\\X^*&I_c\end{smallmatrix}\right)$.
At a singular-value representative $X=(\operatorname{diag}(\sigma_j)\;0)$,
the gradient is real diagonal and this subspace is invariant under the
Hessian.  On its $j$th coordinate the first two derivatives are
$2\sigma_j/(1-\sigma_j^2)$ and
$2(1+\sigma_j^2)/(1-\sigma_j^2)^2$.  Consequently
\[
 \norm{D\Phi(X)}_{D^2\Phi(X),*}^2
       =\sum_{a,j}\frac{2\sigma_j(X_a)^2}{1+\sigma_j(X_a)^2}<R.
\]
The real diagonal section is the cube $(-1,1)^R$, whose sharp barrier
lower bound is $R$ \citep[Proposition 2.3.6]{NN1994}.
This proves the second equality, including coupled competitors.
\end{proof}

Corollary~\ref{cor:spectral-caps} optimizes only the factor count and
dimensions in the whole-row model.
Equation~\eqref{eq:spectral-exact-ledger} gives the exact barrier
parameters of the attaining cones $\mathcal S_G$, not a lower bound for
other, higher-dimensional cones carrying the same rows.
For $k$ identical rows and $g=\lceil k/Q\rceil$, the total dimension and
parameters are
$(M,\nu_{\rm ambient},\nu_{\rm slice})=(k\delta rc+g,kr+g,kr)$.

\subsection{Completion cones and the clique--separator barrier}

Let $G$ be an undirected graph on $n$ vertices; the specified entries
are the diagonal and the positions of edges of $G$.  Over $\R$ or $\C$, let $\mathcal H_G$ be the space of
Hermitian matrices supported on those entries, with the real trace pairing.
The sparse positive-semidefinite cone is
$K_G=\{Y\in\mathcal H_G:Y\succeq0\}$.
Its dual $C_G$ is identified with the partial Hermitian matrices
that admit a positive-semidefinite completion.  This set is closed:
in a convergent sequence of partial matrices all diagonals are
bounded, and every positive-semidefinite completion obeys
$|Z_{ij}|^2\leq Z_{ii}Z_{jj}$.  A convergent subsequence of completions
therefore completes the limit.  Also, $C_G$ is pointed because a matrix
and its negative can both be completed only if all their diagonals, and
hence all entries, vanish.

\begin{proposition}[Classical completion barrier and its exact parameter]
\label{prop:completion-parameter}
For arbitrary graphs $G_1,\ldots,G_l$ of orders $n_1,\ldots,n_l$,
\[
 \vartheta_{\rm opt}\left(\intr\prod_j C_{G_j}\right)=\sum_j n_j.
\]
For one graph an attaining barrier, up to a constant, is
\begin{equation}\label{eq:maxdet-completion-barrier}
 F_G(S)=-\log\det Z(S),
\end{equation}
where $Z(S)$ is the unique maximum-determinant positive-definite completion.
If the graph is chordal, its maximal cliques $\mathcal Q$ and a clique
forest's separators $\mathcal T$, counted with multiplicity, give
\begin{equation}\label{eq:completion-separators}
 F_G(S)=-\sum_{Q\in\mathcal Q}\log\det S_Q
          +\sum_{T\in\mathcal T}\log\det S_T.
\end{equation}
Empty separators contribute zero.
\end{proposition}
\begin{proof}
The restriction of $-\log\det Y$ to $\intr K_G$ is an $n$-logarithmically homogeneous
self-concordant (``$n$-LHSC'') barrier.  Its signed Legendre dual
\[
 F_G(S)=\sup_{Y\in\intr K_G}\{-\ip SY+\log\det Y\}
\]
is another $n$-LHSC barrier \citep{ADV2013}.
Stationarity says that $Y^{-1}$ completes $S$.  Equivalently, stationarity
of the strictly concave completion problem says that $Z(S)^{-1}$ is
sparse.  Boundedness of the set of completions and existence of a
positive-definite completion guarantee existence; strict concavity gives
uniqueness.  Since $Y$ is sparse,
$\ip SY=\tr(Y^{-1}Y)=n$, giving
$F_G=-n-\log\det Z(S)$.
The diagonal section of $C_G$ is precisely $\R_+^n$, and it meets the
interior.  Across a product it is $\R_+^{\sum n_j}$.  The sharp orthant
bound proves the stated lower bound for every coupled barrier.

For a chordal graph, eliminate leaf cliques successively.  If a leaf
clique meets the remainder in separator $T$, Schur complementation
maximizes the determinant by making the residual cross block zero and
multiplies the determinant of the remainder by $\det S_Q/\det S_T$.
Iteration gives \eqref{eq:completion-separators}, the classical
clique--separator formula \citep{GJSW1984,BJL1989}.
The running-intersection identity
$\sum_Q|Q|-\sum_T|T|=n$ gives its degree of logarithmic homogeneity.
Self-concordance of this formula follows from duality, not from a
general rule for subtracting logarithmic barriers.
\end{proof}

\begin{proposition}[Block-star realization]\label{prop:block-star}
Consider a hub clique of size $s\geq1$, and $k$ disjoint leaf cliques
of sizes $p_a\geq1$, each joined completely to the hub.  Put
$\delta=1$ over $\R$ and $\delta=2$ over $\C$, and
$h_\delta(v)=v+\delta v(v-1)/2$.
The completion cone has real dimension and exact ambient parameter
\begin{equation}\label{eq:block-star-ledger}
 D=h_\delta(s)+\sum_a\bigl(\delta sp_a+h_\delta(p_a)\bigr),
 \qquad \nu_{\rm ambient}^{\rm opt}=s+\sum_a p_a.
\end{equation}
Fixing the hub and leaf diagonal blocks to identities gives the product
$\prod_a\{X_a\in\mathbb F^{s\times p_a}:\norm{X_a}_{\rm op}\leq1\}$.
Its restricted completion barrier is exactly
\eqref{eq:spectral-slice-barrier}, with optimal parameter
$\sum_a\min(s,p_a)$, and it represents every full slack row
$1-\ip{X_a}{Z}$ with $Z$ an extreme point of the nuclear-norm unit ball.
\end{proposition}
\begin{proof}
Write the specified blocks as $T$, $U_a$, $W_a$ on the hub, cross edges,
and leaves.  Chordal completion is equivalent to
\[
 \begin{pmatrix}T&U_a\\U_a^*&W_a\end{pmatrix}\succeq0
 \quad\text{for every }a.
\]
The count of independent entries gives $D$, and the graph order gives
the exact parameter by Proposition~\ref{prop:completion-parameter}.
All $k-1$ clique-tree separators are the hub, so the barrier is
\begin{equation}\label{eq:star-barrier}
 -\sum_a\log\det\begin{pmatrix}T&U_a\\U_a^*&W_a\end{pmatrix}
       +(k-1)\log\det T
 =-\log\det T-\sum_a\log\det(W_a-U_a^*T^{-1}U_a).
\end{equation}
At $T=I_s$, $W_a=I_{p_a}$, $U_a=X_a$, the Schur complement gives
exactly the stated product.  Sylvester's determinant identity gives
$\det(I-X_a^*X_a)=\det(I-X_aX_a^*)$, so its barrier is
\eqref{eq:spectral-slice-barrier}, using the smaller order if desired.
For a polar extreme $Z=uv^*$, put $b=(-u,v)$ on the hub and leaf $a$
and zero on all other leaves.  The sparse positive-semidefinite matrix
$bb^*/2$ is a dual factor, and
\[
 \frac12 b^*\begin{pmatrix}I_s&X_a\\X_a^*&I_{p_a}\end{pmatrix}b
       =1-\operatorname{Re}(u^*X_av).
\]
Thus the identity-block section and these dual factors represent
every slack row $1-\ip{X_a}{Z}$.  The rank-one map depends only on $Z$: simultaneous
phase changes in $u,v$ leave $bb^*$ invariant.
\end{proof}

For scalar leaves, $p_a=1$, the construction has
$D=h_\delta(s)+k(\delta s+1)$ and parameter $s+k$.
It removes $\delta k(k-1)/2$ real leaf--leaf coordinates compared with a
dense Hermitian matrix of order $s+k$.  Partitioning $k$ leaves among
$g$ stars gives total dimension $k(\delta s+1)+g h_\delta(s)$ and
parameter $k+gs$.
In the real case, factoring $T$ once and solving against all leaf vectors
evaluates \eqref{eq:star-barrier} in $O(s^3+ks^2)$ dense arithmetic
and stores $O(s^2+ks)$ scalars.  These counts concern barrier evaluation;
a Newton solve with additional constraints has a separate cost.

\subsection{Equality of the restricted barriers and Newton systems}

After fixing the shared scales or the star's diagonal blocks, all
constructions have the same retained coordinates and the same function
$\Phi$ in \eqref{eq:spectral-slice-barrier}.  For one oriented block,
with $M=I-XX^*$, its derivatives are
\[
 \nabla\Phi(X)=2M^{-1}X,\qquad
 D^2\Phi(X)[H]=2M^{-1}H+
       2M^{-1}(HX^*+XH^*)M^{-1}X.
\]
Hence their Hessian metrics and reduced Newton systems are identical:
with the same real affine map $\mathcal A$, right side $b$,
and objective $c$, both systems use
\[
 \begin{pmatrix}D^2\Phi(X)&\mathcal A^*\\\mathcal A&0\end{pmatrix},
 \qquad
 \begin{pmatrix}-D\Phi(X)-\tau c-\mathcal A^*y\\b-\mathcal A X\end{pmatrix}.
\]
Their central paths and reduced condition numbers also coincide.
When orientation requires transposition or adjunction, the identification
is a fixed real orthogonal coordinate change (including conjugation signs
in the complex case).  An oracle for one reduced object therefore yields one for the other
through this known coordinate change.  Building it from either
construction's ambient data is a different task: the ambient
scales, completion coordinates, dual variables, and KKT systems differ.
Neither the smaller ambient parameter nor the smaller sparse coordinate
count alone implies a saving in running time or quantum queries.
